\section{Алгоритмы арифметических операций}
\label{app:algorithms}

В этом приложении представлены блок-схемы алгоритмов арифметических операций. В каждой блок-схеме символом $N$ обозначена разрядность входных операндов.

\def\nodegap{.75cm}

\vspace*{\fill}

\begin{figure}[H]
    \centering
    
    \begin{tikzpicture}[gost flowchart]
    
    % --- Nodes (using anchors and shifts) ---
    \node[terminator, anchor=north] (start) at (0,0) {Начало};
    
    \node[data, anchor=north] (input_ab) at ([yshift=-\nodegap]start.south) {$a, b$};
    
    \node[preparation, anchor=north] (init_ctr) at ([yshift=-\nodegap]input_ab.south) {$ctr = 0$};
    
    \node[preparation, anchor=north] (init_res) at ([yshift=-\nodegap]init_ctr.south) {$part\_res = 0$};
    
    \node[decision, anchor=north] (dec_ctr) at ([yshift=-\nodegap]init_res.south) {$ctr < N-1$};
    
    \node[decision, anchor=north] (dec_bit) at ([yshift=-\nodegap]dec_ctr.south) {$b[ctr] = 1?$};
    
    \node[process, anchor=north] (calc_part) at ([yshift=-\nodegap]dec_bit.south) {$part\_res = part\_res + a \ll ctr$};
    
    \node[process, anchor=north] (inc_ctr) at ([yshift=-\nodegap]calc_part.south) {$ctr = ctr + 1$};
    
    % Right side branch
    \node[process, anchor=west] (assign_res) at ([xshift=2*\nodegap]dec_ctr.east) {$result = part\_res$};
    
    \node[data, anchor=north] (output_res) at ([yshift=-\nodegap]assign_res.south) {$result$};
    
    \node[terminator, anchor=north] (end) at ([yshift=-\nodegap]output_res.south) {Конец};


    % --- Connections (using regular arrow style and explicit anchors) ---

    % Main vertical flow
    \draw[regular arrow] (start.south) -- (input_ab.north);
    \draw[regular arrow] (input_ab.south) -- (init_ctr.north);
    \draw[regular arrow] (init_ctr.south) -- (init_res.north);
    \draw[regular arrow] (init_res.south) -- (dec_ctr.north);
    
    % TRUE branch from ctr < N-1
    \draw[regular arrow] (dec_ctr.south) -- node[right, font=\small] {Yes} (dec_bit.north);
    
    % TRUE branch from b[ctr] == 1
    \draw[regular arrow] (dec_bit.south) -- node[right, font=\small] {Yes} (calc_part.north);
    
    \draw[regular arrow] (calc_part.south) -- (inc_ctr.north);

    % FALSE branch from ctr < N-1 (goes right)
    \draw[regular arrow] (dec_ctr.east) -- node[above, font=\small] {No} (assign_res.west);
    
    % Output to End
    \draw[regular arrow] (assign_res.south) -- (output_res.north);
    \draw[regular arrow] (output_res.south) -- (end.north);

    % FALSE branch from b[ctr] == 1 (routes right, then down, then left to join the main flow)
    % Joins exactly halfway between calc_part and inc_ctr (0.5cm shift since \nodegap is 1cm)
    \draw[regular arrow] (dec_bit.east) 
        -- node[above, font=\small] {No} ([xshift=2.5*\nodegap]dec_bit.east) 
        |- (inc_ctr.east); 

    % Loop back from ctr = ctr + 1 to ctr < N-1
    % Joins exactly halfway between init_res and dec_ctr
    \draw[regular arrow] (inc_ctr.west) 
        -- ([xshift=-2.5*\nodegap]inc_ctr.west) 
        |- (dec_ctr.west);

    \end{tikzpicture}
    \caption{Блок-схема алгоритма вычисления умножения чисел $a$ и $b$.}
    \label{fig:cbrt_algorithm}
\end{figure}

\vspace*{\fill}

\begin{figure}[p]
    \centering
    \begin{tikzpicture}[gost flowchart]
    
    % --- Nodes (using anchors and shifts) ---
    \node[terminator] (start) at (0,0) {Start};
    
    \node[data, anchor=north] (input_x) at ([yshift=-\nodegap]start.south) {$x$};
    
    \node[preparation, anchor=north] (calc_m) at ([yshift=-\nodegap]input_x.south) {$m = 1 \ll (N-2)$};
    
    \node[preparation, anchor=north] (init_y) at ([yshift=-\nodegap]calc_m.south) {$y = 0$};
    
    \node[decision, anchor=north] (dec_m) at ([yshift=-\nodegap]init_y.south) {$m \neq 0$};
    
    \node[process, anchor=north] (calc_b) at ([yshift=-\nodegap]dec_m.south) {$b = y \mid m$};
    
    \node[process, anchor=north] (shift_y) at ([yshift=-\nodegap]calc_b.south) {$y = y \gg 1$};
    
    \node[decision, anchor=north] (dec_x) at ([yshift=-\nodegap]shift_y.south) {$x \geq b$};
    
    \node[process, anchor=north] (sub_x) at ([yshift=-\nodegap]dec_x.south) {$x = x - b$};
    
    \node[process, anchor=north] (or_y) at ([yshift=-\nodegap]sub_x.south) {$y = y \mid m$};
    
    \node[process, anchor=north] (shift_m) at ([yshift=-\nodegap]or_y.south) {$m = m \gg 2$};
    
    % Right side branch
    \node[data, anchor=west] (output_y) at ([xshift=\nodegap]dec_m.east) {$y$};
    
    \node[terminator, anchor=north] (end) at ([yshift=-\nodegap]output_y.south) {End};


    % --- Connections (using your 'join' style and explicit anchors) ---

    % Main vertical flow
    \draw[regular arrow] (start.south) -- (input_x.north);
    \draw[regular arrow] (input_x.south) -- (calc_m.north);
    \draw[regular arrow] (calc_m.south) -- (init_y.north);
    \draw[regular arrow] (init_y.south) -- (dec_m.north);
    
    % TRUE branch from m != 0
    \draw[regular arrow] (dec_m.south) -- node[right, font=\small] {Yes} (calc_b.north);
    
    \draw[regular arrow] (calc_b.south) -- (shift_y.north);
    \draw[regular arrow] (shift_y.south) -- (dec_x.north);
    
    % TRUE branch from x >= b
    \draw[regular arrow] (dec_x.south) -- node[right, font=\small] {Yes} (sub_x.north);
    
    \draw[regular arrow] (sub_x.south) -- (or_y.north);
    \draw[regular arrow] (or_y.south) -- (shift_m.north);

    % FALSE branch from m != 0 (goes right)
    \draw[regular arrow] (dec_m.east) -- node[above, font=\small] {No} (output_y.west);
    
    % Output to End
    \draw[regular arrow] (output_y.south) -- (end.north);

    % FALSE branch from x >= b (routes right, then down to join the main flow)
    \draw[regular arrow] (dec_x.east) 
        -- node[above, font=\small] {No} ([xshift=\nodegap]dec_x.east) 
        |- (shift_m.east);

    % Loop back from m = m >> 2 to m != 0
    \draw[regular arrow] (shift_m.west) 
        -- ([xshift=-\nodegap]shift_m.west) 
        |- (dec_m.west);

\end{tikzpicture}
    \caption{Блок-схема алгоритма вычисления квадратного корня числа $x$.}
    \label{fig:cbrt_algorithm}
\end{figure}

\begin{figure}[p]
    \centering
    \begin{tikzpicture}[gost flowchart]
    
    % --- Nodes (using anchors and shifts) ---
    \node[terminator, anchor=north] (start) at (0,0) {Start};
    
    \node[data, anchor=north] (input_x) at ([yshift=-\nodegap]start.south) {$x$};
    
    \node[preparation, anchor=north] (init_y) at ([yshift=-\nodegap]input_x.south) {$y = 0$};
    
    \node[preparation, anchor=north] (init_s) at ([yshift=-\nodegap]init_y.south) {$s = N - 1$};
    
    \node[decision, anchor=north] (dec_s) at ([yshift=-\nodegap]init_s.south) {$s > -3$};
    
    \node[process, anchor=north] (shift_y) at ([yshift=-\nodegap]dec_s.south) {$y = y \ll 1$};
    
    \node[process, anchor=north] (calc_b) at ([yshift=-\nodegap]shift_y.south) {$b = (3y(y+1)+1) \ll s$};
    
    \node[decision, anchor=north] (dec_x) at ([yshift=-\nodegap]calc_b.south) {$x \geq b$};
    
    \node[process, anchor=north] (sub_x) at ([yshift=-\nodegap]dec_x.south) {$x = x - b$};
    
    \node[process, anchor=north] (inc_y) at ([yshift=-\nodegap]sub_x.south) {$y = y + 1$};
    
    \node[process, anchor=north] (dec_s_update) at ([yshift=-\nodegap]inc_y.south) {$s = s - 3$};
    
    % Right side branch
    \node[data, anchor=west] (output_y) at ([xshift=\nodegap]dec_s.east) {$y$};
    
    \node[terminator, anchor=north] (end) at ([yshift=-\nodegap]output_y.south) {End};


    % --- Connections (using your 'join' style and explicit anchors) ---

    % Main vertical flow
    \draw[regular arrow] (start.south) -- (input_x.north);
    \draw[regular arrow] (input_x.south) -- (init_y.north);
    \draw[regular arrow] (init_y.south) -- (init_s.north);
    \draw[regular arrow] (init_s.south) -- (dec_s.north);
    
    % TRUE branch from S > -3
    \draw[regular arrow] (dec_s.south) -- node[right, font=\small] {Yes} (shift_y.north);
    
    \draw[regular arrow] (shift_y.south) -- (calc_b.north);
    \draw[regular arrow] (calc_b.south) -- (dec_x.north);
    
    % TRUE branch from x >= b
    \draw[regular arrow] (dec_x.south) -- node[right, font=\small] {Yes} (sub_x.north);
    
    \draw[regular arrow] (sub_x.south) -- (inc_y.north);
    \draw[regular arrow] (inc_y.south) -- (dec_s_update.north);

    % FALSE branch from S > -3 (goes right)
    \draw[regular arrow] (dec_s.east) -- node[above, font=\small] {No} (output_y.west);
    
    % Output to End
    \draw[regular arrow] (output_y.south) -- (end.north);

    % FALSE branch from x >= b (routes right, then down to join the main flow)
    \draw[regular arrow] (dec_x.east) 
        -- node[above, font=\small] {No} ([xshift=\nodegap]dec_x.east) 
        |- (dec_s_update.east); % Merges directly into the east anchor of the final process node

    % Loop back from s = s - 3 to S > -3
    \draw[regular arrow] (dec_s_update.west) 
        -- ([xshift=-1.5*\nodegap]dec_s_update.west) 
        |- (dec_s.west);

    \end{tikzpicture}
    \caption{Блок-схема алгоритма вычисления кубического корня числа \(x\).}
    \label{fig:cbrt_algorithm}
\end{figure}

\newpage
\section{Листинги разработанных модулей}
\label{app:listings}

\begin{code}
\caption{Модуль сумматора.}
\label{lst:summer}
\begin{minted}{verilog}
\end{minted}
\end{code}

\begin{code}
\caption{Модуль умножителя.}
\label{lst:multiplier}
\begin{minted}{verilog}
\end{minted}
\end{code}

\begin{code}
\caption{Модуль квадратного корня.}
\label{lst:sqrt}
\begin{minted}{verilog}
\end{minted}
\end{code}

\begin{code}
\caption{Модуль кубического корня.}
\label{lst:cbrt}
\begin{minted}{verilog}
\end{minted}
\end{code}

\begin{code}
\caption{Модуль результата функции.}
\label{lst:func_result}
\begin{minted}{verilog}
\end{minted}
\end{code}

\section{Результаты тестирования разработанных модулей}
\label{app:testing}

\begin{code}
\caption{Вывод разработанного тестового модуля.}
\label{lst:test_output}
\begin{minted}{verilog}
\end{minted}
\end{code}